\documentclass[../main.tex]{subfiles}

\begin{document}

\chapter{Matrices}

Matrices are one of the fundamental components that build linear
algebra. This note discusses the fundamental concepts and properties
of matrices.

\section{Basic Terms and Notations}

As always, let's get started with the definition.

\begin{definition}{}{}
A \textbf{matrix} is a 2-dimensional array composed of entries
called \textbf{elements}. A matrix is said to have \textbf{order}
``$m$ by $n$'' and written as $m \times n$ if there are $m$ rows
and $n$ columns.
\end{definition}

Technically, the elements of a matrix can be anything including
functions. However, the majority of the matrices discussed in
linear algebra will have numeric elements. Let's take a look
at some examples.
\[
A=
\begin{bNiceMatrix}[last-row,last-col]
    1  &  2 &  3 & 4  &
        \Vbrace[horizontal-label]{3}{3 \text{ rows}} \\
    -1 & -2 & -3 & -4 \\
    0  & -1 & 0  & 1  \\
    \Hbrace{4}{4 \text{ columns}} &
\end{bNiceMatrix}
\qquad\quad,\qquad
B =
\begin{bNiceMatrix}[last-row,last-col]
\pi & \Vbrace[horizontal-label]{4}{4 \text{ rows}} \\
e \\
0 \\
1 \\
\Hbrace{1}{1 \text{ column}}
\end{bNiceMatrix}
\]
Here, matrix $A$ has order $3 \times 4$ and matrix $B$ has
order $4 \times 1$. The form of a matrix that has order
$m \times n$ can be generalized as the following.
\[
A=
\begin{bNiceMatrix}[last-row,last-col]
    a_{1,1} & a_{1,2} & a_{1,3} & \cdots & a_{1,n}
        & \Vbrace[horizontal-label]{5}{n \text{ rows}} \\
    a_{2,1} & a_{2,2} & a_{2,3} & \cdots & a_{2,n} \\
    a_{3,1} & a_{3,2} & a_{3,3} & \cdots & a_{3,n} \\
    \vdots  & \vdots  & \vdots  & \ddots & \vdots  \\
    a_{m,1} & a_{m,2} & a_{m,3} & \cdots & a_{m,n} \\
    \Hbrace{5}{m \text{ columns}} &
\end{bNiceMatrix}
\]
It is also worth noting that matrix $A$ is often represented
as $[a_{ij}]$ and the elements as $a_{ij}$ like $a_{12}$ for
$a_{1, 2}$.

\begin{definition}{}{}
An element $a_{ij}$ is said to have row index of $i$ and
column index of $j$.
\end{definition}

As you probably have guessed, we assign special names to some
matrices depending on $m$ and $n$.

\begin{definition}{}{}
\textbf{Square matrices} are matrices with an equal number of rows
and columns. The elements with the same row and column index are
called \textbf{diagonal elements} of the matrix. The collection
of such elements is called \textbf{main diagonal}, or
\textbf{principal diagonal}.
\end{definition}

What if either $m$ or $n$ is equal to $1$? Yes, we do have special
terms for them.

\begin{definition}{}{}
\textbf{Row matrices} are matrices that are composed of only one
row. A \textbf{column matrix} is an array composed of one column.
\end{definition}

If you have a machine learning background, you likely have heard
of the term ``high dimensional arrays or vectors.'' Here, the
term ``dimension'' works the same way.

\begin{definition}{}{}
The \textbf{dimension} of a row or column matrix is the number
of \textbf{components}, or the elements. An
\textbf{$\mathbf{n}$-tuple} is an $n$-dimensional row or column
vector.
\end{definition}

Let's take a look at an example from the previous matrix $B$
reproduced below.

\begin{minipage}{0.2\textwidth}
\[
B =
\begin{bNiceMatrix}[last-row,last-col]
\pi & \Vbrace[horizontal-label]{4}{4 \text{ rows}} \\
e \\
0 \\
1 \\
\Hbrace{1}{1 \text{ column}}
\end{bNiceMatrix}
\]
\end{minipage}
\hfill
\begin{minipage}{0.7\textwidth}
Here, matrix $B$ is a 4-dimensional matrix with four components.
This column matrix can also be called as $4$-tuple.
\end{minipage}

We could also define matrices by their elements.

\begin{definition}{}{}
A \textbf{zero matrix} is a matrix with all its elements as $0$.
\end{definition}

With the fundamental terms in mind, let's discuss the properties
of matrices.

\end{document}


